\section{Conclusion}
\label{sec:conclusion}

% ── Mark scheme: 5 marks ──────────────────────────────────────────────────────

This project designed, implemented, and verified an authoritative Tic-Tac-Toe game server on an STM32F051 microcontroller communicating with a thin Python Tkinter GUI client over USART1 at 115\,200 baud.
The system divides responsibilities strictly: the microcontroller acts as the sole referee, storing the board state, enforcing turn sequence, detecting win/draw outcomes, and updating physical LEDs, while the host GUI handles only user input and graphical display.
The server ensures that rejected requests leave the internal game state invariant before committing any memory changes.

System verification was achieved across two testing tiers:
\begin{itemize}
    \item 48 automated unit and integration tests executing under \texttt{pytest} verified protocol framing, oversized frame recovery, line-ending handling, and client-side error parsing with zero failures.
    \item 26 physical hardware assertions executed over \texttt{COM9} confirmed that the microcontroller's UART responses match the specification byte-for-byte across all compulsory requirements.
    \item End-to-end round-trip latency averaged 4.2\,ms, easily satisfying interactive GUI responsiveness requirements.
\end{itemize}

\subsection{System Limitations}
The current design has three main limitations:
\begin{enumerate}
    \item \textbf{Transport Error Detection}: The application protocol relies on the USB packet checksums provided by the onboard ST-Link Virtual COM Port. The serial payload lacks an application-layer CRC or parity check, meaning byte corruption over a long physical serial line could pass undetected if it matches valid syntax.
    \item \textbf{Single-Client Concurrency}: The server architecture is single-session, communicating with one host client at a time over USART1 without arbitrating between two independent physical serial connections.
    \item \textbf{Volatile Game State}: Game state resides in volatile SRAM (\texttt{g\_board}). A power interruption resets the session to \texttt{IDLE} without persistent match logging in non-volatile memory.
\end{enumerate}

\subsection{Potential Extensions}
Logical future extensions include embedding an optimal Minimax decision algorithm into the microcontroller's Flash memory for standalone single-player matches, configuring a hardware timer (such as \texttt{TIM14}) as a per-turn chess clock, and adding an 8-bit CRC suffix to every serial frame for robust deployment over industrial fieldbuses.
Overall, the project demonstrates that an inexpensive microcontroller can reliably serve as a strict, standalone referee for turn-based systems.
